% INTRODUCTION GÉNÉRALE - start Arabic numbering, fit on 1 page
% ============================================================
\pagenumbering{arabic}
\setcounter{page}{1}

\chapter*{INTRODUCTION GÉNÉRALE}
\addcontentsline{toc}{chapter}{Introduction générale}
\thispagestyle{plain}
\onehalfspacing
\setlength{\parskip}{0.45em}

Depuis plusieurs années, l'agriculture de précision occupe une place croissante dans la production agricole. Elle ne se limite plus à l'usage d'outils mécaniques : elle repose de plus en plus sur des systèmes capables d'observer une culture, d'en interpréter l'état, et d'agir en conséquence. Dans les régions où l'eau est une ressource rare et précieuse, comme c'est le cas ici à Madagascar, le suivi du stress hydrique des plantes devient un enjeu déterminant pour préserver les rendements sans gaspiller une ressource limitée. Or, cette surveillance reste aujourd'hui largement manuelle : elle demande du temps, de la main-d'œuvre, et ne permet qu'une observation ponctuelle d'un champ, difficile à répéter régulièrement à grande échelle.

C'est dans ce contexte qu'intervient le projet Tandem, dont l'objectif est de concevoir un robot suspendu par des câbles au-dessus d'une parcelle, capable de photographier automatiquement les plantes afin d'en surveiller l'état. Un tel robot, appelé robot parallèle à câbles (CDPR), présente l'avantage d'être léger, peu coûteux et adapté à de grands espaces de travail, contrairement à des solutions robotiques classiques plus lourdes. Encore faut-il que ce robot sache où aller et comment s'y déplacer efficacement. C'est précisément l'objet du sous-projet SP2 : définir une méthode de planification et d'optimisation de trajectoire permettant au robot de couvrir une parcelle de manière efficace, tout en respectant les contraintes mécaniques du système et en trouvant le meilleur compromis entre rapidité, économie d'énergie, couverture du terrain et précision des prises de vue.

Cette problématique n'est pas propre au projet Tandem. Dans des domaines très différents, comme la surveillance par drone ou l'inspection robotisée de grandes structures, la question du déplacement optimal d'un système autonome revient systématiquement : qu'il s'agisse d'un robot à câbles, d'un drone ou d'un bras robotisé, la qualité de la trajectoire suivie conditionne directement la qualité des données recueillies et l'efficacité globale du système. Cela montre que la planification de trajectoire n'est pas un simple détail d'implémentation, mais une étape déterminante de la conception de tout système robotique autonome.

\clearpage
\setlength{\parskip}{0.35em}
\onehalfspacing
